% My intro paragraph
When I was younger, about $8$ - $9$ years old, I enjoyed watching shows like
\textbf{Popular Mechanics} and \textbf{National Geographic}. Unfortunately I
regretted watching one of the \textbf{National Geographic} videos for almost $2$
years. The video talked about sinkholes, which terrorized me. Since I was a very
optimistic, naive little kid who didn't think anything was wrong in the world,
this took its toll on me. My bubble of a world shook. The world wasn't all
flowers and roses anymore. I created many dark ideas about how the world may
end. From fire and flames to floods and light beaming from the sky. From that
day on, I imagined every day as the last day. Every day, I would act as if it
was my last day, spending countless hours reading about sinkholes. Every time
I'd close my eyes, I would hear the piercing sound coming from the sinkholes and
visually see it as if I was falling into one. My family forbade me from reading
any more books on the matter. Eventually, like all little kids, I grew out of
this phobia.

% Inline etymological definition of the term apocalypse
The online etymological dictionary defines the term \textbf{Apocalypse} as
\textit{a cataclysmic event}. There are other ways to interpret the term, which
is what I'll be discussion with you.

% Talk about the Quranic End of Times
In the Quranic apocalypse, which is referred to as the \textbf{Final Hour}. It's
described in great detail in the Quran and authentic hadith (which are the
sayings and teachings of the Prophet Muhammad). There are two chapters in the
Quran which are dedicated to talking about this specific event. The first one,
Al-Waqiah, which means \textbf{The Inevitable} and Al-Qiyamah, which means
\textbf{The Resurrection}. He stated that the End of Days won't happen until you
see the $10$ signs.

\begin{enumerate}
  \label{enum:the_10_signs}

  \item The Smoke
  \item The Dajjal (Anit-Christ)
  \item The Beast
  \item The Sun Rising from the West and Setting in the East
  \item The Descent of Jesus (Son of Mary) PBU
  \item The Coming of Gog and Magog
  \item[7-9.] Three Earth Quakes that shall shake this Earth
    \begin{enumerate}
      \label{enum:earth_quakes}
    
      \item One the West
      \item One the East
      \item One in the Arabian Peninsula
    \end{enumerate}
  \item[10.] A Fire that shall appear from Yemen and it Shall gather/force
    people to go towards their place of resurrection
\end{enumerate}

\noindent He continued by saying if you see one of these signs, expected the
next immediately. One comes, expect the rest. For the Dajjal to come, there are
pre-requests. $3$ years before he comes, the first year, a third of the rain
will be held back, the second year, two-thirds of the rain will be held back and
the third year, all the rain will be held back, meaning, when the Dajjal finally
comes down, mankind is suffering from a drought, with him, a river of fire and a
river of water. He enters in a village, amidst the people and he says
\textit{"Do you believe in me? I am your lord."} When they believe, he tells the
sky to rain and rain comes and tells the Earth to produce and it will produce.
He'll travel all around the world, every inch, except for two places. Makkah and
Medina. He will live for $40$ days. The first day will be as long as a year,
second day as long as a month, third day as long as a week, and the rest of the
days normal. After the $40$ days have ended, the Muslims will be praying behind
Al Mahdi, who is the ruler of the Muslims, in Syria in a Mosque called
\textbf{The One With the White Minaret}. While they're praying Duhr, Jesus PBU
will descend. After they finish prayer, Jesus PBU gives glad tidings to all the
people around, informing them of their places in Paradise. They go off to al
Jerusalem to fight the army of the Dajjal, while the Dajjal has come with an
army, not knowing that Jesus is with the Muslims already. They'll find them on
the borders of the Temple of Solomon. When the Dajjal sees Jesus PBU, he runs,
melting while Jesus PBU chases saying \textit{It is written that I owe you on
strike}. Shortly after, he kills the Dajjal. Later on, Al Mahdi will die $10$
years later and within his $10$ years of ruling, he will fill the world with
justice. Shortly after his death, Jesus PBU will die. The sun starts to rise
from where it sets and sets from where it rises. Prophet Muhammad said
\textit{Allah will always accept repentance, unless two things happens, one the
soul reaches the gargling point, and when the sun rises from where it sets.}
From that point, more and more faith is taken away, the Quran is taken away, and
even the Ka'bah is destroyed. At the end, after the other signs take place, a
light breeze will cover the Earth and kill all the remaining Muslims. The Angle,
Israfeel will blow the horn, which shatters everything, brings not just planets,
but the entire universe. Then, he blows it again. After that, Allah will fold
this universe with one hand and the other $6$ universes with his other, and
destroy them. He commands the angle of death to kill all the other angles and to
take his own soul. That's when the Judgment Day begins.

The \textbf{Train to Busan} is an apocalyptic Korean zombie movie. The main plot
of the movie is about a father and his daughter trying to make their way across
\textbf{South Korea}, via a train, when, out of nowhere, a virus starts to
spread among the passengers. The said virus kills the victim and makes them
become a zombie. It focuses on different aspects, which other zombie movies may
not. It criticizes society on many levels, one being, how humans are and how an
event so colossal as an apocalypse can show their true skin. The characters in
the movie make this movie distinct because they're very likable people since
they aren't horrible to other people and selfless. Usually, we see one main
character who survives by sheer luck or by sacrificing everybody to save
himself. In this movie, the clever ones are the people standing in the end if
they decide not to sacrifice themselves to save everybody else.

Initially, I assumed that the term \textbf{Apocalypse} meant the end of the
world. Although that definition is a worldwide accepted meaning, there are other
definitions and interpretations. One interpretation may be \textit{an Apocalypse
brings out the true color in people.} It will help you distinguish your friends
from your enemies and vice versa. Also, an apocalypse may be the end of a
chapter and the beginning of something new like how all the religious scriptures
say. Even in the \textbf{Train to Busan}. The apocalypse may start something new
like a discovery about something from a different world.
